\documentclass[11pt]{article}

\usepackage[a4paper,margin=1in]{geometry}
\usepackage{setspace}
\usepackage{hyperref}

\title{Applications of Artificial Intelligence in Modern Systems}
\author{Generated Example}
\date{\today}

\begin{document}

\maketitle
\onehalfspacing

\section{Introduction}

AI (Artificial Intelligence) has become one of the most influential technologies of the modern era. 
AI refers to computational systems that can perform tasks that normally require human intelligence, 
such as reasoning, perception, and decision-making.

A major subset of AI is ML (Machine Learning), where systems learn patterns from data rather than 
being explicitly programmed. In recent years, ML techniques have been widely deployed in commercial 
systems ranging from recommendation engines to autonomous vehicles.

Another important component of modern computing infrastructure is the API (Application Programming Interface). 
APIs allow different software systems to communicate with each other, enabling AI models to be deployed 
in scalable environments such as cloud platforms.

\section{Data and Infrastructure}

Modern AI systems rely heavily on large-scale data processing. Data is often collected through 
the IoT (Internet of Things), which consists of interconnected devices such as sensors, smart appliances, 
and embedded systems.

These devices communicate over  LAN (Local Area Network) and WAN (Wide Area Network), transmitting 
continuous streams of information to centralized servers.

Once collected, data is stored and managed using a DBMS (Database Management System). 
Large organizations often deploy their DBMS solutions in cloud environments such as AWS (Amazon Web Services)
or GCP (Google Cloud Platform), which provide scalable storage and compute resources.

CPU (Central Processing Unit) and GPU (Graphics Processing Unit) play a crucial role in training AI models. 
While CPUs handle general-purpose tasks, GPUs are optimized for parallel computation, making them ideal 
for deep learning workloads.

\section{Communication and Networking}

AI-powered applications depend on robust communication protocols. The HTTP (HyperText Transfer Protocol) 
is commonly used for web communication, while secure transmission is handled using SSL (Secure Sockets Layer) 
and TLS (Transport Layer Security).

At a lower level, the TCP (Transmission Control Protocol) and IP (Internet Protocol) form the foundation 
of internet communication. The DNS (Domain Name System) translates human-readable domain names into IP addresses, 
allowing users to access services without memorizing numeric identifiers.

These networking technologies ensure that AI systems deployed across distributed environments remain connected, 
responsive, and secure.

\section{Applications of AI}

AI is widely used across multiple industries. In healthcare, AI systems assist doctors in diagnosing diseases 
by analyzing medical images and patient records. In finance, ML (Machine Learning) models are used to detect 
fraudulent transactions in real time.

In transportation, AI powers autonomous driving systems that rely on sensor data processed through IoT networks. 
These systems use APIs to communicate with cloud-based decision engines that continuously improve performance.

In everyday life, AI is embedded in recommendation systems used by streaming platforms, online retailers, and 
social media networks. These systems analyze user behavior to provide personalized content.

\section{Conclusion}

AI, supported by ML, cloud computing, and advanced networking 
technologies, is transforming the way modern systems operate. As infrastructure continues to evolve, technologies 
like APIs, IoT, and GPU-accelerated computing will remain essential for scaling intelligent systems.

A good way to perform these function is using SSH its good bro! Windows Remote Desktop (RDP) is also a good way to perform these functions. The future of networking will require faster interconnections leading to the implementation of hollow-core fibres (HCFs) this can move data at the speed of light. The future of AI will be dependent on the implementation of these technologies.  
\end{document}